\documentclass{article}
\title{Weighted Schnorr Threshold Signatures}
\date{2023-01-19, revised 2026-11-22}
\author{Joseph Wiley Yandle}
\usepackage{amsmath}
\usepackage{amsfonts}
\usepackage{amssymb}
\usepackage{mathtools}
\usepackage[symbol]{footmisc}
\usepackage{caption}
\usepackage{amsthm}

\theoremstyle{plain}
\newtheorem{thm}{Theorem}[section]
\newtheorem{lem}[thm]{Lemma}
\theoremstyle{definition}
\newtheorem{definition}[thm]{Definition}
\theoremstyle{remark}
\newtheorem{rem}[thm]{Remark}

\newcommand\sample{\xleftarrow{\$}}

\renewcommand{\thefootnote}{\fnsymbol{footnote}}

\newcommand\Set[2]{\{\,#1\mid#2\,\}}
\newcommand\underoverset[3]{\underset{#1}{\overset{#2}{#3}}}

\usepackage{geometry}
\geometry{
  a4paper,
  total={170mm,257mm},
  left=20mm,
  top=20mm,
}

\makeatletter
\def\@maketitle{%
  \newpage
  \null
  \vskip 2em%
  \begin{center}%
  \let \footnote \thanks
    {\Huge\bfseries\@title \par}%
    \vskip 1.5em%
    {\large
      \lineskip .5em%
      \begin{tabular}[t]{c}%
        \@author \footnote{xoloki@gmail.com}
      \end{tabular}\par}%
    \vskip 1em%
    {\large \@date}%
  \end{center}%
  \par
  \vskip 1.5em}
\makeatother

\begin{document}
\onecolumn
\maketitle

\begingroup
\leftskip5em
\parskip1em
\rightskip\leftskip
\noindent\textbf{Abstract.} We present Weighted Schnorr Threshold Signatures (WSTS aka WileyProofs), which optimizes the base FROST implementation to reduce bandwidth when FROST parties control multiple keys.  We also prove that WSTS is unforgeable in the random oracle model under the algebraic one-more discrete logarithm assumption, by adapting the proof of Chu, Gerhart, Ruffing and Schr\"oder for FROST3 to parties that hold several keys.

\par
\noindent\textbf{Keywords:} Weighted Aggregate Threshold Signatures; Zero-Knowledge Proofs; FROST
\par
\endgroup

\section{
  Introduction
}

FROST \emph{(Flexible Round-Optimized Schnorr Threshold)} \cite{frost} is a system for making aggregate threshold signatures.  This allows a set of $parties$ to construct a group key, and then sign messages using this key as long as a $threshold$ subset cooperates to form the signature.  This signature is aggregated, so that its size does not depend on the number of signing parties.

The design of FROST assumes that each $party$ controls exactly one key, i.e. the $threshold$ is not weighted.  The naive approach to turning FROST into a weighted threshold scheme involves allowing each $signer$ to control a subset of the $parties$ proportional to its weight.  But this results in inefficiencies in the number of messages required for the protocol, total bandwidth, and computational complexity.

Here we present optimizations on top of vanilla FROST which reduce these inefficiencies when used as a weighted threshold scheme.  We call this scheme WSTS \emph{(Weighted Schnorr Threshold Signatures)}, aka WileyProofs, in honor not only of the author's grandfather, but also the greatest engineer of all time: Wiley E. Coyote.

\subsection{
  Notation
}

While WSTS can be implemented in any group, we will use curve point terminology exclusively.  This will hopefully make it more accessible to engineers wishing to implement it.

Let $G$ be a generator in an elliptic curve group $\mathbb{G}$ of prime order $p$, so that scalars are elements of $\mathbb{Z}_p$.  Let $f$ be a polynomial of a scalar variable with scalar coefficients, and $P$ be a scalar polynomial with curve point coefficients.  Both $f$ and $P$ can be evaluated at a scalar, though the latter will use scalar/point multiplication rather than scalar/scalar.

Let $H$ be a hash function, which maps a series of binary inputs to a fixed size binary digest.  For a hash function to be effective, this mapping must be $one-way$; i.e. it should be impossible to go from a hash digest back to the binary inputs.  The most common hash functions include SHA2, SHA3, and Blake, each of which can create a variety of digest sizes, frequently $256$ or $512$ bits.

\newpage
\onecolumn
\section{
  Background
}

\subsection{
  Threshold Signatures
}

Threshold signatures allow a subset of a group of $N$ signers to collaborate to sign messages.  There are many such protocols, but they all specify a threshold $T$, which is the number of signers necessary to create a valid signature.  

\subsection{
  Aggregate Signatures
}

Traditional approaches to the multisignature problem require that each signer signs separately, which means that the size of the signature grows linearly with the number of required signers.  An aggregate scheme, conversely, condenses these individual signatures into a single group signature.  This allows a large set of signers without increasing the data size necessary to transmit and store the signature.

In a blockchain context, where transaction fees are proportional to the size of the data, this leads to significant savings.  It also prevents the case where a large group signature would exceed the blocksize, thereby artificially limiting the number of signers.

\subsection{
  Aggregate Threshold Signatures
}

By combining the concepts of aggregate and threshold signatures, we arrive at a construction well suited to a blockchain context.  We can have any number of possible signers, and any threshold required to make a valid signature.  This allows signatures which can be used in many more contexts than a traditional multisig construct, while keeping transaction fees low and taking no more blockchain space than a standard single signature.

\subsection{
  Weighted Aggregate Threshold Signatures
}

Aggregate threshold signatures function well in a blockchain context, but the number of messages and associated protocol bandwidth required grow linearly with the total number of parties.  In many use cases, though, not all signers should be given equal weight.  This is particularly relevant to Proof of Stake blockchain systems, where the size of a user's stake gives their vote more weight.

It would be ideal to construct a native weighted threshold scheme, where this size would instead grow proportionally to the number of actual signers, rather than the total number of votes.  

\subsection{
  Polynomial Interpolation
}

Consider a polynomial $f$ of degree $k$:
\begin{align}
  f(x) = a_0 + a_1 x + a_2 x^2 + ... + a_k x^k\nonumber
\end{align}

There are $k+1$ coefficients needed to form a polynomial of degree $k$.  Thus in order to uniquely determine a polynomial of degree $k$, we will need to evaluate it at $k+1$ points, giving us a system of $k+1$ equations in $k+1$ unknowns.  Typically, we evaluate at the points $[1, k+1]$.  This is called polynomial interpolation.

While we can use linear algebra to solve these equations, there is no closed form which can be used algebraically.  Lagrange interpolation \cite{lagrange} is an alternative approach which does provide a closed form.  To interpolate $f(x)$ it uses the Lagrange polynomial $L(x)$ and the Lagrange bases $l_i(x)$:
\begin{align}
  L(x) &= \sum_{i=1}^{k+1} f(x_i) \cdot l_i(x)\\
  l_i(x) &= \prod_{j=1, j \neq i}^{k+1} \frac{x - x_j}{x_i - x_j}
\end{align}

We can see that $l_i(x_i)$ must be $1$, since $\frac{x_i - x_j}{x_i - x_j} = 1$.  For all $m \in [1, k+1] \ni m \neq i$, $l_i(x_m)$ must be $0$, since one of the top terms will be $x_m - x_m$, hence the product must be $0$.  So $\forall i \in [1,k+1];\; L(x_i) = f(x_i)$.

If we only need to evaluate the interpolated polynomial at $0$, we can condense $l_i(0)$ to:
\begin{align}
  l_i(0) &= \prod_{j=1; j \neq i}^{k+1} \frac{x_j}{x_j - x_i}
\end{align}

We call this $\lambda_i$, and use it to quickly evaluate $L(0)$ when given a set $(x_i, f(x_i))$.

\newpage
\onecolumn

\subsection{
  Schnorr Proofs
}\label{section:schnorr}

A Schnorr proof \cite{schnorr} is a zero knowledge proof of ownership of a public key.  Let scalar $x$ be a private key, with point $X = x \cdot G$ the corresponding public key.  To prove knowledge of $x$ given $X$, Schnorr proofs use a $3$-move commitment-challenge-response protocol, i.e. a Sigma protocol.

The prover first commits to a random scalar $v$ with the corresponding point $V$, then sends $V$ to the verifier.  The verifier then returns a random challenge scalar $c$.  The prover then constructs the response $r = v + cx$.  The proof consists of the tuple $(r, c, V)$.

To verify the proof, check that:
\begin{align}
  V = r \cdot G - c \cdot X
\end{align}

This must be true for a valid proof, since:
\begin{align}
  V &= v \cdot G = (r - cx) \cdot G = r \cdot G - cx \cdot G = r \cdot G - c \cdot X\nonumber
\end{align}

To make this protocol non-interactive, use the Fiat-Shamir transform \cite{fiatshamir} to construct $c$ by hashing the initial proof elements $(G, X, V)$:
\begin{align}
  c = H(G, X, V)
\end{align}

\subsection{
  FROST
}

FROST is an aggregate threshold signing scheme.  It allows for a group of $N$ parties to create a distributed group signing key, then some threshold $T$ of them can cooperate to sign a message.  This group signature is an aggregate of the individual party signatures, and functions as a Schnorr proof.

Thus there are three discrete elements required: distributed key generation, gathering party signatures, and forming an aggregate group signature using the party signatures.

\subsubsection{
  Distributed Key Generation (DKG)
}

Like many DKG schemes, FROST employs the technique of polynomial interpolation.  But whereas traditional DKG protocols use a trusted dealer, FROST is a trustless protocol.  To accomplish this, each party acts as a trusted dealer to every other party, and runs a traditional DKG protocol.  The parties then combine their keys trustlessly, and the result is a fully trustless distributed key, of which each party controls an equal share.

Remember that in polynomial interpolation, you need $k+1$ points to interpolate a degree $k$ polynomial.  FROST specifies a threshold $T$, which is the number of signers necessary to sign a message.  Thus we will need a group polynomial of degree $T-1$ so that any group of $T$ signers will be able to interpolate the key, and thereby sign.  This group polynomial $f$ will be the sum of the party polynomials $f_i$; the group private key $x$ will be $f(0)$, and the group public key $Y$ will be $P(0)$, where $P = f \cdot G$.

First, each party is given a sequential ID $i \in [1, N]$, which will be used both to identify it and as an element in the underlying math.  Next, each party $i$ constructs a random private scalar polynomial $f_i$ of degree $T-1$ by sampling its coefficients $a_{i,0}, \ldots, a_{i,T-1} \sample \mathbb{Z}_p$:
\begin{align}\label{equation:frost-polynomial}
  f_i(x) = \sum_{j=0}^{T-1} a_{i,j} \cdot x^j, \qquad P_i(x) = f_i(x) \cdot G.
\end{align}
A public DKG share is constructed from $f_i$ by multiplying each coefficient by the generator $G$, giving the coefficients of $P_i$.  Party $i$ also proves knowledge of the constant term $a_{i,0}$ with a Schnorr proof (Section~\ref{section:schnorr}): it samples $r \sample \mathbb{Z}_p$, sets $R_i = r \cdot G$ and $c_i = H_{\mathrm{pok}}(i, R_i, a_{i,0} \cdot G, \Phi)$, where $\Phi$ is a context string identifying this DKG run, and computes $s_i = r + c_i \cdot a_{i,0}$.  The ID $i$, the public polynomial $P_i$ and the proof $(R_i, s_i)$ are then shared to all other parties, which check that $P_i$ has exactly $T$ coefficients and that $s_i \cdot G = R_i + c_i \cdot P_i(0)$.  The proof rules out rogue-key attacks, in which a party chooses its constant term as a function of the other parties' public values without knowing its discrete logarithm \cite{frost}.

Each party $i$ evaluates its own $f_i$ at each $j \in [1,N]$, and sends a private DKG share to each consisting of $(j, f_i(j))$.  Each receiving party checks each private share for consistency with the public polynomial of the sending party.  This is done for party $i$ by checking that:
\begin{align}\label{equation:frost-dkg-check}
  \forall j \in [1,N] \ni j \neq i;\;f_j(i) \cdot G = P_j(i)
\end{align}

Once the private shares are verified, each party $i$ computes its share $s_i$ of the group secret key by summing the private polynomial evaluations at $i$ from all parties, then uses it to construct its public key $Y_i$:
\begin{align}\label{equation:frost-share}
  s_i = \sum_{j=1}^{N} f_j(i) ;\; Y_i = s_i \cdot G
\end{align}

Each party computes the group public key $Y$ by summing the first coefficient of each public polynomial:
\begin{align}\label{equation:frost-group-key}
  Y = \sum_{j=1}^{N} P_j(0)
\end{align}

At this point DKG is complete, and each party $i$ knows $(s_i, Y_i)$ and the group key $Y$.

\subsubsection{
  Gathering Party Signatures
}

To sign a message $m$, it is necessary to choose the signing set $S$ before any party can sign, with $|S| >= T$; this allows polynomial interpolation to work.  Each party $i \in S$ then constructs a private nonce $(d_i,e_i)$ and corresponding public nonce $(D_i,E_i)$, then sends $(i, D_i, E_i)$ to all other parties in $S$; this forms the set $B$.

Once nonces have been received, each party $i \in S$ computes the binding values $\rho_j$ using $B$ and $m$:
\begin{align}\label{equation:frost-rho}
  \forall j \in S;\;\rho_j = H_\rho(j, m, B)
\end{align}

Next they compute the interpolation value $\lambda_i$:
\begin{align}\label{equation:frost-lambda}
  \lambda_i = \prod_{j \in S; j \neq i}^{} \frac{j}{j - i}
\end{align}

Each party also computes the sum $R$:
\begin{align}\label{equation:frost-R}
  R = \sum_{j \in S}^{} D_j + \rho_j \cdot E_j
\end{align}

Once they have $R$ and message $m$ they can construct the challenge $c$:
\begin{align}\label{equation:frost-c}
  c = H_c(Y, R, m)
\end{align}

Finally, each party $i$ can construct its signature share $z_i$:
\begin{align}\label{equation:frost-z}
  z_i = d_i + \rho_i \cdot e_i + s_i \cdot c \cdot \lambda_i
\end{align}


\subsubsection{
  Aggregating the Group Signature
}

Once the aggregator has gathered the signature shares $z_i$, it must first verify that they are valid by checking that:
\begin{align}\label{equation:frost-share-check}
  z_i \cdot G = D_i + \rho_i \cdot E_i + Y_i \cdot c \cdot \lambda_i
\end{align}

Finally the aggregator can construct the group signature $z$:
\begin{align}\label{equation:frost-z-sum}
  z = \sum_{i \in S}^{} z_i
\end{align}

The full signature consists of $(R, z)$ using $R$ from \eqref{equation:frost-R}.

\subsubsection{
  Verifying the Group Signature
}

To verify the group signature, construct the value $R'$ from the group public key $Y$, challenge $c$, and group signature $z$:
\begin{align}\label{equation:frost-verify}
  R' = z \cdot G - c \cdot Y
\end{align}

If $R = R'$, the proof is valid.

\subsubsection{
  FROST as a Schnorr Proof
}\label{section:frost-schnorr}

How does the FROST algorithm produce a Schnorr proof?  To do so, it should first commit to some randomness $v$ via $V$, produce a challenge that depends on $V$, and then give a response that combines the commitment $v$, challenge $c$, and underlying secret $x$: $r = v + cx$.

This is exactly what FROST does.  The committed values are the nonces $(d_i, e_i)$, shared as $(D_i, E_i)$, and aggregated into the sum $R$, which is used to construct the challenge.  And if we expand the aggregation of the party signatures $z_i$, we can see that the response, the signature $z$, is indeed $v + cx$:
\begin{align}
  z = \sum_{i \in S}^{} z_i = \sum_{i \in S}^{} d_i + \rho_i \cdot e_i + s_i \cdot c \cdot \lambda_i = \sum_{i \in S}^{} (d_i + \rho_i \cdot e_i) + c \cdot \sum_{i \in S}^{}(s_i \cdot \lambda_i)\nonumber
\end{align}

As we know from Lagrange interpolation, $\sum_{i \in S}^{} s_i \cdot \lambda_i$ is the group polynomial evaluated at $0$, which is the group private key $x$.

\newpage
\onecolumn
\section{
  Related Work
}

There are many aggregate threshold signature schemes, some quite similar to FROST.  While FROST is optimized for a reduced number of rounds in the optimistic case of no byzantine actors, there are others which are called $robust$; i.e. they take more rounds in the common case, but can complete even if some signers are byzantine actors.

An excellent example of a robust protocol is Provably Secure Distributed Schnorr Signatures (PSDSS) \cite{psdss}.  The construction features robustness in both the DKG and signing rounds.  At both stages, the protocol will succeed as long as at least $T$ parties complete it honestly.

We did not choose PSDSS as the basis for WSTS for two reasons.  First, rather than gather nonces from signing parties in each signing round, it runs an additional DKG to create an ephemeral round key.  Since DKG is an expensive protocol in the case of a large number of voting slots, it was suboptimal for our use case.  Also, WSTS-style bandwidth reductions did not appear to be viable for PSDSS.

There are also many native weighted threshold signing schemes, but ultimately all were rejected.  Most of them relied on RSA with a trusted setup, and so were considered prima facie unsuitable.

There was one compelling candidate though: An Efficient and Secure Weighted Threshold Signcryption Scheme \cite{eswtss}.  It uses standard elliptic curves combined with a dynamic knapsack sytem, and the Chinese Remainder theorem to set the voting weights.  We rejected it due to its complexity and lack of track record in production systems.  But it remains an ongoing object of research.

The security of FROST itself has been studied in several models.  Komlo and Goldberg \cite{frost} prove a three-round variant, FROST-Interactive, unforgeable under the discrete logarithm assumption and argue heuristically that the two-round protocol inherits the result.  The two-round protocol was later proven secure directly: by Crites, Komlo and Maller in the algebraic group model \cite{ckm21}, by Bellare et al.\ under the one-more discrete logarithm assumption with idealized key generation, together with a hierarchy of unforgeability notions \cite{bckmtz22, btz22}, and by Chu, Gerhart, Ruffing and Schr\"oder under the algebraic one-more discrete logarithm assumption with the key generation included \cite{cgrs23}.  None of these results applies to WSTS as a black box, because an honest WSTS party signs for all of its keys with a single nonce pair; Appendix~\ref{section:security} adapts the last of them.

\section{
  WSTS: Weighted Schnorr Threshold Signatures
}

Our contribution is a set of optimizations on top of vanilla FROST to streamline the protocol in a heavily weighted threshold scenario (e.g. $150$ signers controlling a total of $4000$ keys). In FROST, each party controls a single key, which acts as a single vote.  A naive implementation of a weighted threshold scheme on top of FROST involves giving each signer control of multiple parties, proportional to that signer's weight.  This is functional, but suboptimal.

WSTS optimizes this by having each signer control a single party, but that party now controls multiple shares of the group key.  Crucially, each party still uses only a single polynomial per DKG round and a single nonce per signing round.  This also allows for a single signature per party, rather than per key.  This leads to substantial savings during both DKG and signing.

Whereas FROST contains two primary parameters ($T$ and $N$), WSTS splits the $N$ parameter into two: the number of parties $N_p$, and the number of keys $N_k$.  $T$ is now a threshold of the keys, not the parties.  Each party must still send a DKG private share for each key, but DKG public and signature shares are now per party.

Each party is given not only a party ID $i_p \in [1, N_p]$ as in FROST, but also a set of key IDs $i_k \in [1, N_k]$.  We denote the latter set as $K_{i_p}$.

\subsection{
  Distributed Key Generation (DKG)
}\label{section:wsts-dkg}

Each party $i_p \in [1, N_p]$ begins DKG by generating a random scalar polynomial $f_{i_p}$ of degree $T-1$ as in \eqref{equation:frost-polynomial}.  They then construct DKG public shares by multiplying the coefficients of $f_{i_p}$ by the generator, to obtain $P_{i_p}$, and attach a Schnorr proof of knowledge of the constant term $f_{i_p}(0)$ exactly as in FROST, with the party ID $i_p$ and a context string $\Phi$ bound into the challenge.  The public shares are broadcast, so that every party receives the same set.  A coordinator then announces the set $\mathcal{A} \subseteq [1, N_p]$ of parties whose public shares it accepted; the remaining steps run over $\mathcal{A}$, and a party not in $\mathcal{A}$ abandons the round.  In the common case $\mathcal{A} = [1, N_p]$.

Once the DKG public shares have been distributed, the parties create DKG private shares by evaluating $f_{i_p}$ at all points $j_k \in [1, N_k]$.  These shares consist of a set of tuples $(i_p, (i_k, f_{i_p}(i_k)))$.  Each private share is sent only to the party holding key $i_k$, over a channel that provides confidentiality and authentication.

Once a party receives its DKG public and private shares, it first verifies every public share in $\mathcal{A}$: the proof of knowledge must verify and $P_{j_p}$ must have exactly $T$ coefficients.  It then verifies each private share against the corresponding public share, using a similar formula as before:
\begin{align}\label{equation:wsts-dkg-check}
  \forall j_p \in \mathcal{A}, i_k \in K_{i_p}; f_{j_p}(i_k) \cdot G = P_{j_p}(i_k)
\end{align}

Once the private shares are verified, each party $i_p$ computes its shares $s_{i_k}$ of the group secret key by summing the private polynomial evaluations from all parties for each key $i_k$, giving the group polynomial $f$ evaluated at $i_k$, i.e. one interpolation point. The parties then use each $s_{i_k}$ to construct the corresponding public key $Y_{i_k}$:
\begin{align}\label{equation:wsts-share}
  s_{i_k} = \sum_{j_p \in \mathcal{A}} f_{j_p}(i_k) ;\; Y_{i_k} = s_{i_k} \cdot G
\end{align}

Each party computes the group public key $Y$ by summing the first coefficient of each public polynomial:
\begin{align}\label{equation:wsts-group-key}
  Y = \sum_{j_p \in \mathcal{A}} P_{j_p}(0)
\end{align}

At this point DKG is complete, and each party $i_p$ knows all of its $(s_{i_k}, Y_{i_k})$ and the group key $Y$.

\subsection{
  Gathering Party Signatures
}\label{section:wsts-sign}

As before, in order for polynomial interpolation to work, it is necessary to choose the signing parties $S_p \ni \forall s_p \in S_p, s_p \in [1, N_p]$, and the corresponding $S_k$ which consists of the union of all $i_k \in K_{i_p}\; \forall i_p \in S_p$, before any party can sign, with $|S_k| >= T$.  Each party derives $S_k$ from $S_p$ and the public key assignment, never from key IDs claimed in a signing request; otherwise a dishonest coordinator could make it sign with the wrong interpolation values.  Each party $i_p \in S_p$ then constructs a private nonce $(d_{i_p},e_{i_p})$ and corresponding public nonce $(D_{i_p},E_{i_p})$, then sends $(i_p, D_{i_p},E_{i_p})$ to all other parties with keys in $S_k$, forming set $B$ as before.  A nonce is used for exactly one signature share and then discarded.

Once nonces have been received, each party $i_p \in S_p$ computes the binding values $\rho_{j_p}$ using $B$ and $m$.  Unlike FROST, WSTS also binds the group public key $Y$ into $\rho$:
\begin{align}\label{equation:wsts-rho}
  \forall j_p \in S_p;\;\rho_{j_p} = H_\rho(j_p, Y, m, B)
\end{align}

Next, $\forall i_k \in K_{i_p}$,  they compute the interpolation values $\lambda_{i_k}$:
\begin{align}\label{equation:wsts-lambda}
  \lambda_{i_k} = \prod_{j_k \in S_k; j_k \neq i_k}^{} \frac{j_k}{j_k - i_k}
\end{align}

Each party also computes the sum $R$:
\begin{align}\label{equation:wsts-R}
  R = \sum_{j_p \in S_p}^{} D_{j_p} + \rho_{j_p} \cdot E_{j_p}
\end{align}

Once they have $R$ and message $m$ they can construct the challenge $c$:
\begin{align}\label{equation:wsts-c}
  c = H_c(Y, R, m)
\end{align}

Finally, each party $i_p$ can construct its signature share $z_{i_p}$:
\begin{align}\label{equation:wsts-z}
  z_{i_p} = d_{i_p} + \rho_{i_p} \cdot e_{i_p} + \sum_{j_k \in K_{i_p}}^{} s_{j_k} \cdot c \cdot \lambda_{j_k}
\end{align}


\subsection{
  Aggregating the Group Signature
}

Once the aggregator has gathered the signature shares $z_{i_p}$, it must first verify that they are valid by checking that:
\begin{align}\label{equation:wsts-share-check}
  z_{i_p} \cdot G = D_{i_p} + \rho_{i_p} \cdot E_{i_p} + \sum_{j_k \in K_{i_p}}^{} Y_{j_k} \cdot c \cdot \lambda_{j_k}
\end{align}

Finally the aggregator can construct the group signature $z$:
\begin{align}\label{equation:wsts-z-sum}
  z = \sum_{i_p \in S_p}^{} z_{i_p}
\end{align}

The full signature consists of $(R, z)$ using $R$ from \eqref{equation:wsts-R}.

\subsection{
  Verifying the Group Signature
}

To verify the group signature, construct the value $R'$ from the group public key $Y$, challenge $c$, and group signature $z$:
\begin{align}\label{equation:wsts-verify}
  R' = z \cdot G - c \cdot Y
\end{align}

If $R = R'$, the proof is valid.

\subsubsection{
  WSTS as a Schnorr Proof
}\label{section:wsts-schnorr}

WSTS constructs a Schnorr proof in the same way as FROST, with the same commitment value and challenge.  The response, the signature $z$, likewise expands to $v + cx$:
\begin{align}
  z = \sum_{i_p \in S_p}^{} z_{i_p} = \sum_{i_p \in S_p}^{} d_{i_p} + \rho_{i_p} \cdot e_{i_p} + \sum_{j_k \in K_{i_p}}^{} s_{j_k} \cdot c \cdot \lambda_{j_k} = \sum_{i_p \in S_p}^{} (d_{i_p} + \rho_{i_p} \cdot e_{i_p}) + c \cdot \sum_{j_k \in S_k}^{} (s_{j_k} \cdot \lambda_{j_k})\nonumber
\end{align}

As we know from Lagrange interpolation, $\sum_{j_k \in S_k}^{} s_{j_k} \cdot \lambda_{j_k}$ is the group polynomial evaluated at $0$, which is the group private key $x$.


\subsection{
  Hash Functions
}\label{section:hashes}

WSTS uses three hash functions, each mapping to a scalar and each under its own domain separation tag, so that they can be modelled as independent random oracles:
\begin{itemize}
  \item $H_{\mathrm{pok}}(i_p, R, A, \Phi)$, the challenge of the DKG proof of knowledge.  The reference implementation computes SHA-256 over the tag \texttt{WSTS/polynomial-constant}, the party ID, the compressed points $R$ and $A$, and the context $\Phi$, which is the DKG round identifier.
  \item $H_\rho(j_p, Y, m, B)$, the binding value.  The implementation first hashes the commitment list $B$ to a scalar with the \texttt{expand\_message\_xmd} construction of RFC 9380 under the tag \texttt{WSTS/group\_commitment}, then expands the party ID, the compressed group key, the message and that digest under the tag \texttt{WSTS/binding}.
  \item $H_c(Y, R, m)$, the Schnorr challenge.  The implementation uses the BIP340 tagged hash \texttt{BIP0340/challenge} \cite{bip340} over the x-coordinates of $R$ and $Y$ and the message, so that the group signature can be checked by any BIP340 verifier.
\end{itemize}
Nonces are derived as in RFC 9591 \cite{rfc9591}, by hashing fresh randomness together with a secret held by the party, so that a faulty random number generator alone cannot produce a zero or repeated nonce.

\subsection{
  Signature Variants
}\label{section:variants}

The protocol as described produces a signature $(R, z)$ on the full point $R$ that verifies as $z \cdot G = R + c \cdot Y$; we call this the plain variant.  The reference implementation also produces BIP340 signatures \cite{bip340}, which are verified against the x-coordinates of $R$ and of the group key, optionally tweaked as in BIP341 for taproot.  To do so, each party negates its nonce contribution $d_{i_p} + \rho_{i_p} \cdot e_{i_p}$ when $R$ has an odd $y$-coordinate, and negates its key contribution $c \cdot \sum_{j_k} s_{j_k} \cdot \lambda_{j_k}$ when the group key has an odd $y$-coordinate, or, in the taproot case, when exactly one of the group key and the tweaked key does; the aggregator then adds $c$ times the tweak.  The security result of Appendix~\ref{section:security} is stated for the plain variant.

\newpage
\onecolumn
\section{
  Performance
}

The WSTS Rust reference implementation \cite{wstsimpl} contains vanilla FROST in the v1 module, and WSTS in the v2 module.  This crate contains criterion benchmarks for both v1 and v2.

The results are quite staggering.  When testing DKG with $4$ signers and $20$ keys, v2 outperforms v1 by an order of magnitude: $53 ms$ per DKG round for v1 compared to $8.3 ms$ for v2.  Similar speedups were noted for party signing and group signing.

In the table, $N_s$ will refer to the number of signers, which is $N$ in FROST and $N_p$ in WSTS.  $N_k$ is the number of keys, and $T$ is the threshold.

\begin{center}
  \captionof{table}{FROST vs WSTS (runtime)}
  \begin{tabular}{ c|c|c|c|c|c|c|c|c } 
    \multicolumn{3}{c|}{} & \multicolumn{2}{|c|}{DKG}& \multicolumn{2}{|c|}{Party Sign} & \multicolumn{2}{|c}{Group Sign} \\
    \hline
    \hline
    $N_s$ & $N_k$ & $T$ & v1 & v2 & v1 & v2 & v1 & v2 \\ 
    \hline
    \hline
    4 & 20 & 13 & 53ms & 8.3ms & 18ms & 0.80ms & 2.1ms & 0.85ms \\ 
    4 & 40 & 13 & 200ms & 15ms & 120ms & 2.0ms & 6.5ms & 2.2ms \\ 
    4 & 40 & 26 & 420ms & 36ms & 120ms & 2.0ms & 6.6ms & 2.3ms \\ 
    4 & 60 & 26 & 1000ms & 59ms & 380ms & 4.0ms & 13ms & 4.6ms \\ 
    4 & 60 & 40 & 1900ms & 120ms & 380ms & 4.0ms & 13ms & 4.6ms \\ 
    4 & 80 & 40 & 3400ms & 160ms & 880ms & 6.7ms & 23ms & 7.6ms \\ 
    4 & 80 & 53 & 4900ms & 240ms & 880ms & 6.7ms & 23ms & 7.6ms \\ 
    4 & 100 & 53 & 7900ms & 300ms & 1700ms & 10ms & 35ms & 12ms \\ 
    4 & 100 & 66 & 10000ms & 400ms & 1700ms & 10ms & 35ms & 12ms \\ 
  \end{tabular}
\end{center}


\section{
  Robustness
}

As discussed in Related Work, there are some aggregated threshold schemes which are robust, in that they can complete even in the presence of byzantine actors, as long as there exists a threshold set which completes the protocol honestly.  These were rejected for a variety of reasons, as discussed.

However, it is possible to develop a meta protocol on top of a non-robust scheme which provides robustness.  It simply requires running a non-robust scheme repeatedly, removing the bad actors as they are exposed.  Eventually the meta protocol will converge to success, if a threshold set of honest actors exists, or to failure if it does not.  Here we will explore two such meta protocols: FIRE and ROAST.

\subsection{
  FIRE
}

In addition to WSTS, here we also present FROST Interactive Robustness Extension (FIRE).  This is a straightforward extension of FROST into a meta protocol which provides robustness.  It works with both FROST and WSTS, with minor differences in the protocol.

A FIRE signing round takes place after a DKG round has established a group key.  Each round is labeled with an integer ID $i$, and consists of a series of sessions $j$, where $j$ starts at $0$ for every round.

Let $A$ be the set of active parties.  When a signing round begins, $A$ will contain all parties who successfully completed DKG.  We will remove parties from $A$ as the protocol runs, which will allow us to excise byzantine actors.

So when round $i$ begins, we request nonces for session $j = 0$ from $A$.  Once we get $k = T$ nonces for $i,j$, the corresponding parties will be selected for the signing set $S_j$.  We then request signature shares from all parties in $S_j$.  If all sign correctly, and we obtain a valid signature, round $i$ is complete.  We must set a timeout for the signing portion of the session, or else byzantine actors can slow the protocol indefinitely.

Parties who do not return signatures, or return invalid signatures, are removed from $A$.  Parties who did not return a nonce for round $j$ are also removed.  We then we begin session $j+1$.  All members of $A$ are then requested to give nonces for session $j+1$.  The same process happens as before; when we get $k = T$ nonces we form $S_{j+1}$ and again request signature shares.

As before, if the aggregate signature is valid, round $i$ is complete.  Otherwise, we begin a new session $j+2$, and continue in the same vein.  Since each session will either complete or remove some parties from $A$, this algorithm will eventually terminate.  When $|A| < T$, signing round $i$ has failed.

When running a FIRE round with WSTS vs FROST, the only difference is that we must count the number of keys controlled by each party who responds with a nonce in every session.  Only when the sum of keys controlled by the given nonces equals or exceeds $T$ do we form $S_j$ and begin the signing portion of the session.  


\subsection{
  ROAST
}

ROAST\cite{roast} is a wrapper around FROST and other threshold signature schemes which provides robustness and asynchronicity.  It operates in a very similar manner to FIRE, but the asynchonicity allows for faster completion if there exists a set of honest and responsive signers.

ROAST differs from FIRE by allowing session $j+1$ to run in parallel with session $j$.  To do so, it is necessary to keep track of two sets: $R$ is the set of responsive signers, and $M$ of known malicious signers.  Once a party becomes a member of set $M$, all messages from that party will be ignored, and if the governance model includes sanctions on bad actors, those in $M$ will be nominated for such.

The ROAST paper describes the protocol in terms of events and responses, which makes an initial read somewhat challenging.  We will describe it here more linearly.  So as with FIRE, signing round $i$ begins session $j = 0$ with the coordinator asking all signers for nonces.  $R$ and $M$ are initially the null set.  Signers who receive a request to begin $i,0$ respond with a nonce $n_0$ only; this will be different $\forall j \neq 0$.

As the coordinator receives responses, the responding signers are placed into $R$, and their nonces are placed in $\rho_j$.  Once $|R| = T$, the coordinator sends $(\rho_j, R)$ to the signers to request signature shares.  Signers respond with signature $z_j$, and also a nonce $n_{j+1}$.  Piggybacking the next session's nonce with the current session's signature share improves the responsiveness of the protocol.

Crucially, when the coordinator sends $(\rho_j, R)$ to signers, those signers in $R$ are removed.  As they respond, the signers are placed back into $R$ if their signature shares are valid; if they are not valid, those signers are placed into $M$ and declared malicious.  Signers who respond with nonces for the current round while the signing session continues are placed into $R$ with those who provided valid signatures.

As signature shares and nonces are received, the coordinator does two things.  It first checks to see if it has received a full set of valid signature shares; if this happens, and the group signature is valid, then round $i$ is complete.  At the same time, the size of $R$ is examined.  When $|R| = T$, session $j+1$ begins.  Crucially, this happens in parallel to session $j$.

As the sessions continue for each signing round, it is clear that if there exists a threshold set of honest and responsive signers, the protocol will complete succesfully.  Every session will either succeed or add malicious actors to $M$.  If $|M| > (N-T)$, then the signing round has failed.  Likewise, we will never have more than $N-T$ sessions for any round.  Though the paper does not describe the use of timeouts, without them we will never be able to declare that a round is over, since a byzantine set of actors can simply stall any session where one of the set is a participant.  But since the session timeouts run in parallel, the speedup even in the byzantine failure case is significant.

As with FIRE, we can use either FROST or WSTS as the signature aggregation protocol, with the only difference being how we count the nonces when determining if we have a threshold set of possible signers.

\section{
  Acknowledgements
}

The author would like to thank Alie Slade for initial work on the WSTS algorithm.

\begin{thebibliography}{99}

\bibitem{frost}
  Chelsea Komlo, Ian Goldberg
  \emph{FROST: Flexible Round-Optimized Schnorr Threshold Signatures} 2020.12.22.
  \texttt{https://eprint.iacr.org/2020/852.pdf}

\bibitem{lagrange}
  \emph{Lagrange polynomial}.
  \texttt{https://en.wikipedia.org/wiki/Lagrange\%5Fpolynomial}

\bibitem{schnorr}
  \emph{Schnorr signature}.
  \texttt{https://en.wikipedia.org/wiki/Schnorr\%5Fsignature}

\bibitem{fiatshamir}
  \emph{Fiat Shamir heuristic}.
  \texttt{https://en.wikipedia.org/wiki/Fiat\%2DShamir\%5Fheuristic}

\bibitem{psdss}
  D.R. Stinson, R. Strobl
  \emph{Provably Secure Distributed Schnorr Signatures and a (t,n) Threshold Scheme for Implicit Certificates} ACISP 2001. Lecture Notes in Computer Science, vol 2119 2001.01.23.
  \texttt{https://doi.org/10.1007/3-540-47719-5\_33}

\bibitem{eswtss}
  Chien-Hua Tsai
  \emph{An Efficient and Secure Weighted Threshold Signcryption Scheme} Journal of Internet Technology, vol. 20, no. 5 , pp. 1523-1534, Sep. 2019
  \texttt{https://jit.ndhu.edu.tw/article/view/2135}

\bibitem{wstsimpl}
  Joey Yandle
  \emph{WSTS Reference Implementation}
  \texttt{https://github.com/xoloki/wsts}

\bibitem{roast}
  Tim Ruffing, Viktoria Ronge, Elliott Jin, Jonas Schneider-Bensch, Dominique Schroder
  \emph{ROAST: Robust Asynchronous Schnorr Threshold Signatures} 
  \texttt{https://eprint.iacr.org/2022/550.pdf}

\bibitem{ckm21}
  Elizabeth Crites, Chelsea Komlo, Mary Maller
  \emph{How to Prove Schnorr Assuming Schnorr: Security of Multi- and Threshold Signatures} 2021.
  \texttt{https://eprint.iacr.org/2021/1375}

\bibitem{bckmtz22}
  Mihir Bellare, Elizabeth Crites, Chelsea Komlo, Mary Maller, Stefano Tessaro, Chenzhi Zhu
  \emph{Better than Advertised Security for Non-Interactive Threshold Signatures} CRYPTO 2022.

\bibitem{btz22}
  Mihir Bellare, Stefano Tessaro, Chenzhi Zhu
  \emph{Stronger Security for Non-Interactive Threshold Signatures: BLS and FROST} 2022.
  \texttt{https://eprint.iacr.org/2022/833}

\bibitem{cgrs23}
  Hien Chu, Paul Gerhart, Tim Ruffing, Dominique Schr\"oder
  \emph{Practical Schnorr Threshold Signatures Without the Algebraic Group Model} CRYPTO 2023.
  \texttt{https://eprint.iacr.org/2023/899}

\bibitem{nrs21}
  Jonas Nick, Tim Ruffing, Yannick Seurin
  \emph{MuSig2: Simple Two-Round Schnorr Multi-Signatures} CRYPTO 2021.
  \texttt{https://eprint.iacr.org/2020/1261}

\bibitem{bn06}
  Mihir Bellare, Gregory Neven
  \emph{Multi-Signatures in the Plain Public-Key Model and a General Forking Lemma} ACM CCS 2006.

\bibitem{bcj08}
  Ali Bagherzandi, Jung Hee Cheon, Stanis\l{}aw Jarecki
  \emph{Multisignatures Secure under the Discrete Logarithm Assumption and a Generalized Forking Lemma} ACM CCS 2008.

\bibitem{rfc9591}
  Deirdre Connolly, Chelsea Komlo, Ian Goldberg, Christopher A. Wood
  \emph{RFC 9591: The Flexible Round-Optimized Schnorr Threshold (FROST) Protocol for Two-Round Schnorr Signatures} 2024.
  \texttt{https://www.rfc-editor.org/rfc/rfc9591}

\bibitem{bip340}
  Pieter Wuille, Jonas Nick, Tim Ruffing
  \emph{BIP 340: Schnorr Signatures for secp256k1}.
  \texttt{https://github.com/bitcoin/bips/blob/master/bip-0340.mediawiki}

\end{thebibliography}

\appendix
\newpage
\onecolumn

\section{
  Security
}\label{section:security}

\subsection{
  Correctness
}

In Sections~\ref{section:wsts-schnorr} and~\ref{section:frost-schnorr}, we show that the WSTS and FROST group signatures are Schnorr signatures under the group public key: $(R, z)$ satisfies $z \cdot G = R + c \cdot Y$ with $c = H_c(Y, R, m)$.  This is sufficient to show correctness.

\subsection{
  Overview
}

The remainder of this appendix proves that WSTS is unforgeable.  A natural first attempt is to reduce WSTS to FROST as a black box, by viewing each key as a FROST participant.  That attempt fails, for a reason worth recording: a WSTS party signs for all of its keys with a \emph{single} nonce pair, whereas the FROST participants it would map to each carry an independent nonce pair and an independent binding value.  No reduction can turn the latter into the former without the discrete logarithms of the nonces, so the FROST theorem cannot be invoked for any configuration in which an honest party holds more than one key, which is every configuration WSTS exists for.

Instead we prove unforgeability directly, by adapting the proof of Chu, Gerhart, Ruffing and Schr\"oder \cite{cgrs23} for Olaf (FROST3 with a Pedersen key generation with proofs of possession).  Their proof uses only three properties of the way the secret key is shared among the signers: for an authorized signing set the honest and corrupted additive shares sum to the secret key; each honest signer's additive share has a public key computable from the key generation transcript; and the simulator can sample the corrupted shares consistently with the group key.  WSTS has all three, with the honest party's additive share $\sum_{k \in K_i} \lambda_k s_k$ in place of FROST's $\lambda_i s_i$.  Two further differences from Olaf are handled explicitly: WSTS uses a per-party binding value as in FROST1 \cite{btz22} rather than a single one, and WSTS parties are corrupted as units carrying several keys, so the usual convention of padding the corrupted set to exactly $T-1$ participants is replaced by padding the corrupted \emph{key} set.

The result is stated for the plain variant of Section~\ref{section:variants}.  Like \cite{cgrs23} it is in the random oracle model under the algebraic one-more discrete logarithm assumption, a falsifiable weakening of the one-more discrete logarithm assumption, and it does not use the algebraic group model.

\subsection{
  Model
}\label{section:model}

\paragraph{Parameters.}  Throughout, $\mathbb{G}$ is a group of prime order $p$ with generator $G$, where $p$ is a $\lambda$-bit prime, so $2^{\lambda-1} \le p < 2^{\lambda}$.  A WSTS instance has $N_p$ parties and $N_k$ keys with threshold $T$; the \emph{key assignment} $K$ partitions $[1, N_k]$ into the sets $K_i$, $i \in [1, N_p]$.  To lighten the notation of the body we write $i$ for a party and $k$ for a key, and for a set of parties $S$ we write $w(S) = \sum_{i \in S} |K_i|$ for its weight and $\mathrm{keys}(S) = \bigcup_{i \in S} K_i$ for its key set.  For a set of keys $S_k$ and $k \in S_k$, $\lambda_k^{S_k}$ is the Lagrange coefficient of \eqref{equation:wsts-lambda}.

\paragraph{Corruption.}  An adversary statically corrupts a set $\mathcal{C} \subset [1, N_p]$ of parties with $w(\mathcal{C}) < T$ before the protocol begins, and controls them completely.  It also plays the coordinator and the signature aggregator.  The remaining parties $\mathcal{H} = [1, N_p] \setminus \mathcal{C}$ are honest and follow the protocol.  Note that $\mathcal{C}$ may contain more than $T-1$ parties, and $\mathcal{H}$ may contain only one.

\paragraph{Channels.}  During DKG, public shares and the coordinator's messages are sent over an authenticated broadcast channel: every honest party receives the same messages, from honest and corrupted parties alike.  This is the assumption made for FROST by Komlo and Goldberg \cite{frost}, who assume that ``participants do not provide different commitment values to a subset of participants''; the protocol of \cite{cgrs23} instead rebuilds this agreement with an explicit equality check.  Private shares are sent over secure point-to-point channels, which the reference implementation realizes by encrypting them under a Diffie-Hellman key.  Signing messages travel over a network controlled by the adversary.

\paragraph{Random oracles.}  The three hash functions $H_{\mathrm{pok}}$, $H_\rho$ and $H_c$ of Section~\ref{section:hashes} are modelled as independent random oracles with range $\mathbb{Z}_p$.  Each is queried by the adversary at most $q_h$ times.

\begin{definition}[Algebraic one-more discrete logarithm \cite{nrs21}]\label{definition:aomdl}
  An algorithm $\mathcal{A}$ for the AOMDL problem in $\mathbb{G}$ is given access to a challenge oracle $\mathcal{O}_{\mathrm{ch}}$, which on each call returns a fresh uniformly random $X_j \in \mathbb{G}$, and a discrete logarithm oracle $\mathcal{O}_{\mathrm{dl}}$, which takes a group element $Z$ together with an explicit representation $Z = \alpha \cdot G + \sum_j \beta_j \cdot X_j$ in terms of the challenges obtained so far and returns $\log_G Z$.  $\mathcal{A}$ wins if it returns the discrete logarithms of all $c$ challenges it requested after having made fewer than $c$ queries to $\mathcal{O}_{\mathrm{dl}}$.  The AOMDL assumption states that every efficient $\mathcal{A}$ wins with negligible probability $\mathrm{Adv}^{\mathrm{aomdl}}_{\mathcal{A}, \mathbb{G}}$.
\end{definition}

The representation requirement is what makes AOMDL falsifiable; it is weaker than OMDL, in which $\mathcal{O}_{\mathrm{dl}}$ takes a bare group element.  Our reduction always knows a representation for the elements it queries.

We use the forking lemma of Bellare and Neven \cite{bn06} and the multi-forking lemma of Bagherzandi, Cheon and Jarecki \cite{bcj08}, in the formulation of \cite[Lemmas~1 and~2]{cgrs23}.  Both are stated for a randomized algorithm $\mathcal{A}$ that takes an input $\mathit{inp}$, a sequence $h_1, \ldots, h_q$ of elements of $\mathbb{Z}_p$ (the random oracle replies it will hand out), and random coins, and returns either $\bot$ or an index into the sequence together with a side output.  Its accepting probability $\mathrm{acc}_{\mathcal{A}}$ is the probability that it does not return $\bot$.

\begin{lem}[Forking lemma \cite{bn06}]\label{lemma:bn}
  Let $\mathrm{Fork}_{\mathcal{A}}$ run $\mathcal{A}$ once, and if it returns an index $f$ run it a second time with the same input and coins but with $h_f, \ldots, h_q$ replaced by fresh values, succeeding if the second run returns the same index $f$ and $h_f \ne h'_f$.  Then $\mathrm{Fork}_{\mathcal{A}}$ succeeds with probability $\mathrm{frk} \ge \mathrm{acc}_{\mathcal{A}} \cdot (\mathrm{acc}_{\mathcal{A}} / q - 1/p)$.
\end{lem}

\begin{lem}[Multi-forking lemma \cite{bcj08}]\label{lemma:bcj}
  Let $\mathcal{A}$ return a nonempty set $F$ of indices instead of a single one.  Let $\mathrm{MFork}_{\mathcal{A}}$ run $\mathcal{A}$ once and then, for each $f \in F$, repeatedly rerun it with fresh $h_f, \ldots, h_q$, at most $k_{\max} = |F| \cdot 8q / \mathrm{acc}_{\mathcal{A}} \cdot \ln(8|F| / \mathrm{acc}_{\mathcal{A}})$ times, until a run returns a set containing $f$ with $h'_f \ne h_f$.  If $p > 8q |F| / \mathrm{acc}_{\mathcal{A}}$ then, with probability at least $\mathrm{acc}_{\mathcal{A}} / 8$, $\mathrm{MFork}_{\mathcal{A}}$ obtains such a run for every $f \in F$.
\end{lem}

\subsection{
  Unforgeability
}\label{section:unforgeability}

We define unforgeability through a game between a challenger, who runs the honest parties, and the adversary, following the game of \cite[Fig.~2]{cgrs23} with WSTS's key assignment added.  The adversary interacts with the honest parties through three oracles.

\begin{definition}[Unforgeability game]\label{definition:game}
  The game is parameterized by $(N_p, N_k, T, K, \mathcal{C})$ with $w(\mathcal{C}) < T$ and a DKG context string $\mathit{ctx}$.  It initializes $\mathit{Signed} \leftarrow \emptyset$, and gives the adversary the following oracles.
  \begin{itemize}
    \item \textbf{KeyGen}$(i)$, for honest $i$ not yet started: runs party $i$'s side of the DKG of Section~\ref{section:wsts-dkg}, concurrently with the other oracle calls.  The adversary supplies the messages of the corrupted parties and of the coordinator, subject to the channel assumptions above; in particular the coordinator's accepted set $\mathcal{A} \subseteq [1, N_p]$ is broadcast.  If party $i$ completes DKG, the oracle records its group public key as $Y_i$ and its shares $s_k$, $k \in K_i$; otherwise $Y_i = \bot$.
    \item \textbf{Nonce}$(i, j)$, for honest $i$ with $Y_i \ne \bot$ and a fresh session index $j$: party $i$ samples $(d_{ij}, e_{ij})$ and returns $(D_{ij}, E_{ij})$.
    \item \textbf{Sign}$(i, j, S_p, B, m)$, for honest $i$ with $Y_i \ne \bot$, an open session $j$ of party $i$, a set of parties $S_p \ni i$, a list $B$ of one commitment pair per party of $S_p$, and a message $m$: party $i$ closes session $j$, sets $S_k = \mathrm{keys}(S_p)$, computes $\rho_l = H_\rho(l, Y_i, m, B)$ for $l \in S_p$, $R = \sum_{l \in S_p} D_l + \rho_l \cdot E_l$ with $(D_l, E_l)$ taken from $B$, $c = H_c(Y_i, R, m)$, and returns
      \begin{align}\label{equation:game-share}
        z_{ij} = d_{ij} + \rho_i \cdot e_{ij} + c \cdot \sum_{k \in K_i} \lambda_k^{S_k} \cdot s_k.
      \end{align}
      The oracle adds $m$ to $\mathit{Signed}$.
  \end{itemize}
  The adversary wins if it outputs $(i^*, m^*, (R^*, z^*))$ with $Y_{i^*} \ne \bot$, $z^* \cdot G = R^* + H_c(Y_{i^*}, R^*, m^*) \cdot Y_{i^*}$, and $m^* \notin \mathit{Signed}$.  The scheme is $(\tau, q_s, q_h, \epsilon)$-unforgeable if every adversary running in time $\tau$ and making at most $q_s$ queries to each of \textbf{Nonce} and \textbf{Sign} and $q_h$ to each random oracle wins with probability at most $\epsilon$.
\end{definition}

\begin{rem}\label{remark:game}
  Three features of the game deserve comment.  First, the honest party derives the key set $S_k$ from the public key assignment, never from the request; Section~\ref{section:wsts-sign} requires the same of implementations.  Second, the honest party does not check that $w(S_p) \ge T$, that $S_p \subseteq \mathcal{A}$, or that $B$ contains its own commitment; the proof does not need these checks, so omitting them from the game only strengthens the result.  Third, each nonce pair is used in at most one \textbf{Sign} query, which the reference implementation enforces by discarding the nonce after signing.  The winning condition is the weakest in the hierarchy of \cite{btz22}: a message counts as signed as soon as one honest party has produced a share for it.
\end{rem}

\subsection{
  Theorem
}

\begin{thm}[Unforgeability of WSTS]\label{theorem:wsts}
  Fix $(N_p, N_k, T, K, \mathcal{C})$ with $w(\mathcal{C}) < T$.  Let $\mathcal{F}$ be an adversary against the unforgeability of WSTS (plain variant) that runs in time $\tau$, makes at most $q_s$ queries to each of \textbf{Nonce} and \textbf{Sign} and at most $q_h$ queries to each random oracle, and wins with probability $\epsilon$.  Let
  \begin{align}
    q = (N_p + 2)(q_h + q_s) + 1
  \end{align}
  and assume $2^{\lambda - 1} \ge 32 N_p q^2 / \epsilon^2$.  Then there is an algorithm $\mathcal{D}^*$ that solves AOMDL in $\mathbb{G}$ with probability
  \begin{align}\label{equation:theorem-advantage}
    \mathrm{Adv}^{\mathrm{aomdl}}_{\mathcal{D}^*, \mathbb{G}} \;\ge\; \frac{\epsilon^2}{8q} - \frac{N_p + 4q^2}{2^{\lambda}}
  \end{align}
  in expected time at most
  \begin{align}\label{equation:theorem-time}
    \tau' \;\le\; \frac{32 N_p^2 q^2}{\epsilon^2} \cdot \ln\!\left(\frac{16 N_p q}{\epsilon^2}\right) \cdot (\tau_{\mathrm{game}} + \tau),
  \end{align}
  where $\tau_{\mathrm{game}}$ is the time to run the challenger of the game, excluding $\mathcal{F}$.
\end{thm}

The structure of the bound is that of \cite[Theorem~1]{cgrs23}, with the threshold $t$ of that statement replaced by $N_p$: the multi-forking step extracts one proof of knowledge per corrupted party in the accepted set, of which there are at most $N_p - 1$, and the hash-query count $q$ grows by a factor $N_p$ because a single binding query of WSTS fixes $N_p$ binding values at once.

\subsection{
  Proof
}

\begin{proof}
  We construct four algorithms in turn.  $\mathcal{B}$ wraps $\mathcal{F}$ and simulates the game to it, embedding a discrete logarithm challenge $X^*$ in the constant term of every honest party's polynomial and answering signing queries with the oracle $\mathcal{O}_{\mathrm{dl}}$.  $\mathcal{C}$ forks $\mathcal{B}$ on the random oracle query of the forgery and extracts the group secret key $x$.  $\mathcal{D}$ multi-forks $\mathcal{C}$ on the proof-of-knowledge queries of the corrupted parties, extracts their constant terms, and from those and $x$ recovers $x^* = \log_G X^*$, after which it solves every nonce challenge.  $\mathcal{D}^*$ is $\mathcal{D}$ with all but one execution of $\mathcal{C}$ cut off after DKG, which brings the number of $\mathcal{O}_{\mathrm{dl}}$ queries below the number of challenges.

  \paragraph{Notation.}  Let $\mathcal{C}_k = \mathrm{keys}(\mathcal{C})$ be the corrupted keys, so $|\mathcal{C}_k| = w(\mathcal{C}) \le T - 1$.  Fix once and for all a set $E$ of honest keys with $|\mathcal{C}_k \cup E| = T - 1$; this is possible because $N_k \ge T - 1$.  For $k \in [1, N_k]$ and the $T$ interpolation points $\{0\} \cup \mathcal{C}_k \cup E$, write $L_0(k)$ and $L_j(k)$, $j \in \mathcal{C}_k \cup E$, for the Lagrange basis polynomials evaluated at $k$, so that every polynomial $f$ of degree at most $T-1$ satisfies
  \begin{align}\label{equation:proof-interpolation}
    f(k) = L_0(k) \cdot f(0) + \sum_{j \in \mathcal{C}_k \cup E} L_j(k) \cdot f(j).
  \end{align}

  \paragraph{Algorithm $\mathcal{B}$.}  $\mathcal{B}$ takes as input a challenge $X^* \in \mathbb{G}$, challenges $U_{i,1}, \ldots, U_{i,2q_s} \in \mathbb{G}$ for every honest party $i$, a sequence $(h^{\mathrm{pok}}_1, \ldots, h^{\mathrm{pok}}_q)$ for $H_{\mathrm{pok}}$, and a sequence $(h_1, \ldots, h_q)$ shared by $H_\rho$ and $H_c$, all elements of $\mathbb{Z}_p$.  It has access to $\mathcal{O}_{\mathrm{dl}}$.  It keeps tables $\mathcal{T}_{\mathrm{pok}}$, $\mathcal{T}_\rho$ and $\mathcal{T}_c$ of programmed oracle values, counters into the two sequences, a set $\mathcal{Q}$ of the $\mathcal{O}_{\mathrm{dl}}$ queries it makes, and a flag $\mathit{Ev}_1$.  It runs $\mathcal{F}$ with random coins and answers its queries as follows.

  \emph{Random oracle queries.}  A query $H_{\mathrm{pok}}(i, R, A, \mathit{ctx})$ not yet in $\mathcal{T}_{\mathrm{pok}}$ is assigned the next value of the first sequence.  A query $H_\rho(l, Y, m, B)$ with $B$ a well-formed commitment list over a set of parties $S_p \ni l$ is answered in a batch: for every $l' \in S_p$ in increasing order, $\mathcal{T}_\rho(l', Y, m, B)$ is assigned the next value of the second sequence if not yet assigned; then $\mathcal{B}$ computes $R = \sum_{l' \in S_p} D_{l'} + \rho_{l'} \cdot E_{l'}$ from $B$ and assigns $\mathcal{T}_c(Y, R, m)$ the next value if not yet assigned.  Any other $H_\rho$ query, and any $H_c$ query, is assigned the next value of the second sequence if not yet assigned.  The batch ensures that the challenge of every $(Y, m, B)$ is fixed no later than its binding values, which is what the forking argument below relies on; it is the per-party analogue of the internal $H_{\mathrm{sig}}$ query made inside $H_{\mathrm{non}}$ in \cite{cgrs23}.

  \emph{Key generation.}  For every honest party $i$, $\mathcal{B}$ simulates the DKG as follows, which is the simulator of \cite[Fig.~7]{cgrs23} applied to every honest party rather than to one chosen at random.
  \begin{enumerate}
    \item Sample $u_i \in \mathbb{Z}_p$ and set $A_{i,0} = X^* + u_i \cdot G$.  Sample $\tilde{f}_i(j) \in \mathbb{Z}_p$ for every $j \in \mathcal{C}_k \cup E$; these are the values of party $i$'s polynomial at the corrupted and padding keys.  Compute the remaining coefficients $A_{i,1}, \ldots, A_{i,T-1}$ in the exponent from $A_{i,0}$ and $\{\tilde{f}_i(j) \cdot G\}$ by inverting the Vandermonde system of \eqref{equation:proof-interpolation}; this fixes $P_i$ as the commitment to the unique polynomial $f_i$ of degree $T-1$ with $f_i(0) = x^* + u_i$ and $f_i(j) = \tilde{f}_i(j)$ for $j \in \mathcal{C}_k \cup E$.
    \item Simulate the proof of knowledge of $a_{i,0}$: sample $\tilde{s}_i, \tilde{c}_i \in \mathbb{Z}_p$, set $\tilde{R}_i = \tilde{s}_i \cdot G - \tilde{c}_i \cdot A_{i,0}$, and program $\mathcal{T}_{\mathrm{pok}}(i, \tilde{R}_i, A_{i,0}, \mathit{ctx}) \leftarrow \tilde{c}_i$.  If that entry was already assigned, set $\mathit{Ev}_1$ and stop.
    \item Broadcast $(i, P_i, \tilde{R}_i, \tilde{s}_i)$, and send $\tilde{f}_i(j)$ to the corrupted party holding key $j$ for every $j \in \mathcal{C}_k$.  Shares for honest keys are not computed: their recipients are simulated by $\mathcal{B}$ and never need them, and over the secure channels the adversary observes only that a message was sent.
    \item Receive the public shares, proofs and private shares of the other parties and the coordinator's accepted set $\mathcal{A}$, and perform party $i$'s checks exactly as the protocol prescribes, except that the private shares from honest parties are deemed valid.  If any check fails or $i \notin \mathcal{A}$, party $i$ outputs $Y_i = \bot$; otherwise $Y_i = \sum_{j \in \mathcal{A}} A_{j,0}$.
  \end{enumerate}
  By the broadcast assumption every honest party that completes has the same $\mathcal{A}$ and the same key, which we call $Y$.  Writing $x = \log_G Y$, $a_{j,0}$ for the constant term of corrupted party $j$, and $h = |\mathcal{A} \cap \mathcal{H}| \ge 1$,
  \begin{align}\label{equation:proof-key}
    x = h \cdot x^* + \sum_{i \in \mathcal{A} \cap \mathcal{H}} u_i + \sum_{j \in \mathcal{A} \cap \mathcal{C}} a_{j,0}.
  \end{align}
  The share of an honest key $k$ is $s_k = \sum_{j \in \mathcal{A}} f_j(k)$.  The contributions $f_j(k)$ of corrupted parties $j$ were sent to the honest holder of $k$, which $\mathcal{B}$ simulates, and were verified against $P_j$.  The contribution of an honest party $l$ follows from \eqref{equation:proof-interpolation} applied to $f_l$, whose values at $\mathcal{C}_k \cup E$ were chosen by $\mathcal{B}$.  Hence the share of every honest key is an affine function of $x^*$ whose coefficients $\mathcal{B}$ knows as scalars,
  \begin{align}\label{equation:proof-share-affine}
    s_k = \gamma_k \cdot x^* + \delta_k, \qquad \gamma_k = h \cdot L_0(k), \qquad \delta_k = \sum_{l \in \mathcal{A} \cap \mathcal{H}} \Big( L_0(k) \cdot u_l + \sum_{j \in \mathcal{C}_k \cup E} L_j(k) \cdot \tilde{f}_l(j) \Big) + \sum_{j \in \mathcal{A} \cap \mathcal{C}} f_j(k).
  \end{align}
  Consequently, for every honest party $i$ and key set $S_k \supseteq K_i$, the additive share $x_i(S_k) = \sum_{k \in K_i} \lambda^{S_k}_k s_k$ that party $i$ signs with is $\gamma_i(S_k) \cdot x^* + \delta_i(S_k)$ with $\gamma_i(S_k) = \sum_{k \in K_i} \lambda^{S_k}_k \gamma_k$ and $\delta_i(S_k) = \sum_{k \in K_i} \lambda^{S_k}_k \delta_k$, and its public key $X_i(S_k) = x_i(S_k) \cdot G$ has the representation $\gamma_i(S_k) \cdot X^* + \delta_i(S_k) \cdot G$.  Shares of corrupted keys are never needed.

  The simulated DKG is distributed exactly as the real one.  The constant terms $x^* + u_i$ of the honest parties are independent and uniform because the $u_i$ are; each $f_i$ is uniform among polynomials of degree $T-1$ with that constant term because its values at the $T-1$ points $\mathcal{C}_k \cup E$ are independent and uniform; and the simulated proof of knowledge is distributed as a real one unless $\mathit{Ev}_1$ occurs.  Since $\tilde{R}_i$ is uniform in $\mathbb{G}$ and at most $q$ entries of $\mathcal{T}_{\mathrm{pok}}$ are ever assigned,
  \begin{align}\label{equation:proof-ev1}
    \Pr[\mathit{Ev}_1] \le \frac{N_p \cdot q}{p} \le \frac{N_p \cdot q}{2^{\lambda - 1}}.
  \end{align}

  \emph{Nonce queries.}  On the $j$-th \textbf{Nonce} query to honest party $i$, $\mathcal{B}$ returns $(D_{ij}, E_{ij}) = (U_{i,2j-1}, U_{i,2j})$.

  \emph{Signing queries.}  On \textbf{Sign}$(i, j, S_p, B, m)$ for an open session $j$, $\mathcal{B}$ sets $S_k = \mathrm{keys}(S_p)$, obtains $\rho_i = \mathcal{T}_\rho(i, Y, m, B)$ and $c = \mathcal{T}_c(Y, R, m)$ through the batch procedure, and queries
  \begin{align}\label{equation:proof-dl-query}
    z_{ij} = \mathcal{O}_{\mathrm{dl}}\big(U_{i,2j-1} + \rho_i \cdot U_{i,2j} + c \cdot X_i(S_k)\big),
  \end{align}
  supplying the representation $(c \cdot \delta_i(S_k); \; 1 \text{ at } U_{i,2j-1}, \; \rho_i \text{ at } U_{i,2j}, \; c \cdot \gamma_i(S_k) \text{ at } X^*)$.  It records $(i, j, z_{ij}, \rho_i, c, S_k)$ in $\mathcal{Q}$, closes the session, adds $m$ to $\mathit{Signed}$ and returns $z_{ij}$.  By \eqref{equation:game-share} and \eqref{equation:proof-share-affine} this is exactly the share an honest party would return.

  \emph{Output.}  When $\mathcal{F}$ halts with $(i^*, m^*, (R^*, z^*))$, $\mathcal{B}$ returns $\bot$ if $Y_{i^*} = \bot$, if the signature does not verify under $Y$ with $c^* = \mathcal{T}_c(Y, R^*, m^*)$ (assigning it if necessary), if $m^* \in \mathit{Signed}$, or if any proof of knowledge of a corrupted party in $\mathcal{A}$ fails to verify (it cannot, since an honest party completed).  Otherwise let $f_{\mathrm{sig}}$ be the index of $c^*$ in the second sequence, and let $F$ be the set of indices in the first sequence at which $\mathcal{T}_{\mathrm{pok}}(j, \tilde{R}_j, A_{j,0}, \mathit{ctx})$ was assigned for $j \in \mathcal{A} \cap \mathcal{C}$; these are distinct because the party index is part of the query.  $\mathcal{B}$ returns $f_{\mathrm{sig}}$ together with $\big(z^*, F, \{\tilde{s}_j\}_{j \in \mathcal{A} \cap \mathcal{C}}, \{u_i\}_{i \in \mathcal{A} \cap \mathcal{H}}, \mathcal{Q}\big)$.

  \emph{Accounting.}  $H_{\mathrm{pok}}$ is assigned at most $q_h + N_p \le q$ values.  The second sequence is consumed by at most $q_h$ adversarial $H_\rho$ queries and $q_s$ signing queries, each assigning at most $N_p + 1$ values, plus $q_h$ adversarial $H_c$ queries and one final assignment, so at most $(N_p + 1)(q_h + q_s) + q_h + 1 \le q$ values.  Unless $\mathit{Ev}_1$ occurs the simulation is perfect, so
  \begin{align}\label{equation:proof-accb}
    \mathrm{acc}_{\mathcal{B}} \ge \epsilon - \frac{N_p \cdot q}{2^{\lambda - 1}}.
  \end{align}

  \paragraph{Algorithm $\mathcal{C}$.}  $\mathcal{C}$ takes the challenges and the first sequence as input, samples the second sequence, and runs $\mathrm{Fork}_{\mathcal{B}}$ on the index $f_{\mathrm{sig}}$ of the second sequence, relaying $\mathcal{O}_{\mathrm{dl}}$ queries.  If forking succeeds it obtains two forgeries from executions that are identical up to the assignment $\mathcal{T}_c(Y, R^*, m^*) \leftarrow h_{f_{\mathrm{sig}}}$, so both are on the same $(Y, R^*, m^*)$ with challenges $c^* = h_{f_{\mathrm{sig}}} \ne h'_{f_{\mathrm{sig}}} = c'^*$, and
  \begin{align}\label{equation:proof-x}
    x = \frac{z^* - z'^*}{c^* - c'^*}.
  \end{align}
  $\mathcal{C}$ returns $F$ and the proof responses $\{\tilde{s}_j\}$ from the first execution, together with $\theta = (x, \{u_i\}, \mathcal{Q}, \mathcal{Q}')$, where $\mathcal{Q}'$ is the query set of the second execution.  By Lemma~\ref{lemma:bn} and \eqref{equation:proof-accb}, using $\epsilon \le 1$,
  \begin{align}\label{equation:proof-accc}
    \mathrm{acc}_{\mathcal{C}} \;\ge\; \frac{\mathrm{acc}_{\mathcal{B}}^2}{q} - \frac{\mathrm{acc}_{\mathcal{B}}}{p} \;\ge\; \frac{\epsilon^2}{q} - \frac{2N_p + 1}{2^{\lambda - 1}} \;\ge\; \frac{\epsilon^2}{2q},
  \end{align}
  where the last step uses the hypothesis $2^{\lambda-1} \ge 32 N_p q^2 / \epsilon^2$.

  \paragraph{Algorithm $\mathcal{D}$.}  $\mathcal{D}$ requests $X^*$ and the $2 q_s |\mathcal{H}|$ nonce challenges from $\mathcal{O}_{\mathrm{ch}}$, samples the first sequence, and runs $\mathrm{MFork}_{\mathcal{C}}$ on the index set $F$, relaying and caching $\mathcal{O}_{\mathrm{dl}}$ queries.  The hypothesis of Lemma~\ref{lemma:bcj} holds because $|F| \le N_p - 1$ and $\mathrm{acc}_{\mathcal{C}} \ge \epsilon^2 / 2q$ give $8q|F| / \mathrm{acc}_{\mathcal{C}} \le 16 N_p q^2 / \epsilon^2 < 2^{\lambda - 1} \le p$.  If it succeeds, $\mathcal{D}$ holds for every $j \in \mathcal{A} \cap \mathcal{C}$ two valid proofs $(\tilde{R}_j, \tilde{s}_j)$ and $(\tilde{R}_j, \tilde{s}'_j)$ on the same $A_{j,0}$ with distinct challenges, from which $a_{j,0} = (\tilde{s}_j - \tilde{s}'_j) / (\tilde{c}_j - \tilde{c}'_j)$.  By \eqref{equation:proof-key},
  \begin{align}\label{equation:proof-xstar}
    x^* = \frac{1}{h}\Big(x - \sum_{i \in \mathcal{A} \cap \mathcal{H}} u_i - \sum_{j \in \mathcal{A} \cap \mathcal{C}} a_{j,0}\Big).
  \end{align}

  It remains to solve the nonce challenges.  With $x^*$ and the $a_{j,0}$ known, every $\delta_k$ and hence every $x_i(S_k)$ is now a known scalar.  For each honest $i$ and session $j \le q_s$, $\mathcal{D}$ looks in $\mathcal{Q}$ and $\mathcal{Q}'$ for the tuples $(i, j, z, \rho, c, S_k)$ and $(i, j, z', \rho', c', S'_k)$ of the two executions of $\mathcal{B}$ inside the primary execution of $\mathcal{C}$.  Let $\mathit{Ev}_2$ be the event that two of the at most $2q$ values of the second sequence drawn across the two executions coincide; since they are uniform in $\mathbb{Z}_p$,
  \begin{align}\label{equation:proof-ev2}
    \Pr[\mathit{Ev}_2] \le \frac{(2q)^2}{2p} \le \frac{4q^2}{2^{\lambda}}.
  \end{align}
  We claim that unless $\mathit{Ev}_2$ occurs, whenever both tuples exist,
  \begin{align}\label{equation:proof-star}
    \rho = \rho' \;\Longrightarrow\; c = c' \;\wedge\; S_k = S'_k.
  \end{align}
  Suppose $\rho = \rho'$.  All values are distinct, so $\rho$ and $\rho'$ were handed out at the same index $f = f'$ of the second sequence, and since the values at indices $f_{\mathrm{sig}}$ and later are resampled in the second execution, $f < f_{\mathrm{sig}}$.  The executions are identical before $f_{\mathrm{sig}}$, so the assignments at index $f$ have the same key $(i, Y, m, B)$ in both, hence $m = m'$, $B = B'$ and $S_k = S'_k$.  That assignment happened inside the batch for $(Y, m, B)$, which also fixed $\mathcal{T}_c(Y, R, m)$ at some index $f_c$, either earlier, if it was already assigned, or within the batch's consecutive block of indices.  If $f_c < f_{\mathrm{sig}}$ then $c = c'$.  The forgery index $f_{\mathrm{sig}}$ is a $\mathcal{T}_c$ assignment, and the only $\mathcal{T}_c$ assignment in the block is $f_c$, so $f_c > f_{\mathrm{sig}} > f$ is impossible; and $f_c = f_{\mathrm{sig}}$ would mean $(Y, R^*, m^*) = (Y, R, m)$, so $m^* = m$ is a message with a signing query and $\mathcal{B}$ would have returned $\bot$.  This proves \eqref{equation:proof-star}.

  Now, exactly as in \cite{cgrs23}: if the two tuples are identical, the cached oracle answered once and $\mathcal{D}$ makes a second query with fresh $(\rho', c', S'_k)$, $\rho' \ne \rho$; if one or both are missing, $\mathcal{D}$ makes the missing queries likewise; and if both exist and differ, then $\rho \ne \rho'$ by \eqref{equation:proof-star}, since equal $(\rho, c, S_k)$ give equal oracle inputs and hence equal answers.  In every case $\mathcal{D}$ has two equations
  \begin{align}
    u_{i,2j-1} + \rho \cdot u_{i,2j} = z - c \cdot x_i(S_k), \qquad u_{i,2j-1} + \rho' \cdot u_{i,2j} = z' - c' \cdot x_i(S'_k)
  \end{align}
  with $\rho \ne \rho'$, and solves them for the discrete logarithms of $U_{i,2j-1}$ and $U_{i,2j}$.  By Lemma~\ref{lemma:bcj}, \eqref{equation:proof-accc} and \eqref{equation:proof-ev2},
  \begin{align}\label{equation:proof-accd}
    \mathrm{acc}_{\mathcal{D}} \;\ge\; \frac{\mathrm{acc}_{\mathcal{C}}}{8} - \Pr[\mathit{Ev}_2] \;\ge\; \frac{\epsilon^2}{8q} - \frac{2N_p + 1}{2^{\lambda + 2}} - \frac{4q^2}{2^{\lambda}} \;\ge\; \frac{\epsilon^2}{8q} - \frac{N_p + 4q^2}{2^{\lambda}}.
  \end{align}

  \paragraph{Algorithm $\mathcal{D}^*$.}  $\mathcal{D}$ solves all of its challenges but makes too many $\mathcal{O}_{\mathrm{dl}}$ queries, because every execution of $\mathcal{C}$ started by $\mathrm{MFork}$ answers signing queries.  $\mathcal{D}^*$ is $\mathcal{D}$ with one change: in every execution of $\mathcal{C}$ other than the first, both executions of $\mathcal{B}$ are aborted at the moment the first honest party completes DKG, or when $\mathcal{F}$ halts if no honest party ever does.  At that moment the accepted set, the corrupted parties' proofs of knowledge and the outcome of their verification are all determined, so the output that $\mathrm{MFork}$ needs from the aborted execution, namely the index set and the proof responses $\tilde{s}'_j$, can be read off its state.  An honest party answers no signing query before it has completed DKG, so the aborted executions make no $\mathcal{O}_{\mathrm{dl}}$ queries at all.  Everything else $\mathcal{D}$ uses comes from the first execution of $\mathcal{C}$, which is unchanged.  Hence $\mathrm{acc}_{\mathcal{D}^*} = \mathrm{acc}_{\mathcal{D}}$.

  $\mathcal{D}^*$ requests $2 q_s |\mathcal{H}| + 1$ challenges.  Its $\mathcal{O}_{\mathrm{dl}}$ queries all come from the two executions of $\mathcal{B}$ in the first execution of $\mathcal{C}$, each of which makes at most $q_s |\mathcal{H}|$ queries counting those emulated afterwards, so it makes at most $2 q_s |\mathcal{H}|$ queries and wins the AOMDL game whenever it does not return $\bot$.  This gives \eqref{equation:theorem-advantage}.  For the running time, $\mathrm{MFork}_{\mathcal{C}}$ runs at most $|F| \cdot k_{\max} \le 8 N_p^2 q / \mathrm{acc}_{\mathcal{C}} \cdot \ln(8 N_p / \mathrm{acc}_{\mathcal{C}})$ executions of $\mathcal{C}$, each of which runs $\mathcal{B}$ twice, and $\mathcal{B}$ runs in time about $\tau_{\mathrm{game}} + \tau$; substituting $\mathrm{acc}_{\mathcal{C}} \ge \epsilon^2 / 2q$ gives \eqref{equation:theorem-time}.
\end{proof}

\begin{rem}[Where the weighting enters]
  The weighting of WSTS is confined to three places in the proof: the padding set $E$, which replaces the FROST convention of corrupting exactly $T-1$ participants; the coefficients $\gamma_i(S_k)$ and $\delta_i(S_k)$ in \eqref{equation:proof-share-affine}, which sum over a party's keys where FROST has a single Lagrange term; and the fact that the key set $S_k$ in \eqref{equation:proof-star} is determined by the party set $S_p$ through the public key assignment.  Nothing else distinguishes an honest WSTS party from an honest FROST1 participant whose share happens to be $x_i(S_k)$, which is why the single nonce pair per party costs nothing in the argument.
\end{rem}

\begin{rem}[Variants and stronger notions]
  The theorem is stated for the plain variant.  The BIP340 variants of Section~\ref{section:variants} negate the honest party's nonce contribution when $R$ has odd $y$ and the sign of $c \cdot x_i(S_k)$ depending on the parity of the (tweaked) group key; the simulation of \eqref{equation:proof-dl-query} accommodates both by adjusting the representation, and the challenge hash over x-only coordinates adds a sign case to \eqref{equation:proof-x}, since the two forgeries may have $R'^* = -R^*$.  We have not carried this bookkeeping through and leave it for a future revision.  Likewise, the stronger unforgeability notions of \cite{btz22}, in which a message counts as signed only once enough honest weight has signed it, require counting key weight rather than honest signers; FROST1 achieves the third tier of that hierarchy \cite[Theorem~5.4]{btz22} and we expect the same argument to carry over.
\end{rem}

\end{document}
